\section{Application : exemple du poison}

Dans cette partie, nous allons présenter un exemple concret d'application du travail présenté précédemment. Cette section est organisée sous forme de pipeline, nous partons d'un problème défini en PDDL que nous convertissons en LTS, nous ajoutons une évaluation experte avant d'appliquer différentes configurations de la fonction de coût pour en analyser les nuances. L'exemple défini correspond à une situation dans laquelle se trouve un verre qui contient à l'origine de l'eau. On considère que différents événements peuvent se produire pour altérer la potabilité de ce dernier et entraîner ou non la mort de la personne qui le boit. Les événements sont évalués selon une norme que l'on laisse floue, mais qui correspond approximativement à une norme du bien dans le sens où un événement est jugé comme peu normal s'il est associé à une nuisance  (comme empoisonner le verre qui peut entrainer la mort) et comme très normal s'il entraîne un bénéfice (verser un antidote dans un verre empoisonné permet de sauver une personne). La modélisation de la situation et de l'évaluation est entièrement subjective et pourrait être complexifié, mais nous souhaitons garder un exemple simple pour plus de lisibilité. 

\subsection{De PDDL vers un LTS}

La définition du 

\begin{equation*}
    \begin{aligned}
        \mathds{F} & = \{VerreOK, \textit{VerreEmpoisonné},VerreVide,VerreAntidote,Mort,Vivant\}\\ 
        \mathds{E} & = \{empoisonner,antidote,boire,vider\}\\
        \mathds{A} & = \{ Vivant,Mort\}\\
        \mathds{T}& = \mathds{N} \\
        & pre(boire)  = VO \vee VA \vee VE \vee VV \quad pre(antidote) = VE \\
        & pre(empoisonner) = VO \quad pre(vider) = VO\vee VE \\
        \mathds{S} & = \left\{
        \begin{aligned}
        VO     &= VO     \wedge \neg VE \wedge \neg VV \wedge \neg VA \wedge \neg \text{Mort} \wedge \neg \text{Vivant}, \\
        VE     &= \neg VO \wedge VE     \wedge \neg VV \wedge \neg VA \wedge \neg \text{Mort} \wedge \neg \text{Vivant}, \\
        VV     &= \neg VO \wedge \neg VE \wedge VV     \wedge \neg VA \wedge \neg \text{Mort} \wedge \neg \text{Vivant}, \\
        VA     &= \neg VO \wedge \neg VE \wedge \neg VV \wedge VA     \wedge \neg \text{Mort} \wedge \neg \text{Vivant}, \\
        \text{Mort}   &= \neg VO \wedge \neg VE \wedge \neg VV \wedge \neg VA \wedge \text{Mort}   \wedge \neg \text{Vivant}, \\
        \text{Vivant} &= \neg VO \wedge \neg VE \wedge \neg VV \wedge \neg VA \wedge \neg \text{Mort} \wedge \text{Vivant}
        \end{aligned}
        \right\}\\
        \tau & = \left\{
        \begin{aligned}
        (VE,vider,VV),(VE,boire,Mort),(VE,antidote,VA),(Mort,\emptyset,Mort)\\
        (VV,\emptyset,Mort),(VA,boire,Vivant),(Vivant,\emptyset,Vivant),(VO, boire,vivant)\\
        (VO, boire,vivant),(VO,empoisonner,VE),(VO,vider,VV),\\
        \end{aligned}
        \right\} \\
        expert\_judge_\tau &= \left\{
        \begin{aligned}
        VO &: \{boire: 1,\ vider: 0.5,\ emp: 0.1\}, \\
        VE &: \{antidote: 1,\ boire: 0.1,\ vider: 0.5\}, \\
        VV &: \{\emptyset: 0.5\}, \\
        VA &: \{boire: 1\}, \\
        \text{Mort} &: \{\emptyset: 1\}, \\
        \text{Vivant} &: \{\emptyset: 1\}
        \end{aligned}
        \right\}
    \end{aligned}
\end{equation*}



\begin{align}
    \scalebox{1.2}{
        \xymatrix@=2em@C=5em{
            &    *++[o][F-]{Vivant} \ar @(r,u) []_{\emptyset}^1& \\
            &   &    *++[o][F-]{VA} \ar [lu]_{boire}^1  \\
            *++[o][F-]{VO}  \ar [ruu]^{boire}_1 \ar [r]^{empoisonner}_{0.1} \ar[rdd]^{vider}_{0.5} &*++[o][F-]{VE} \ar[ru]^{antidote}_1 \ar[rd]^{boire}_{0.1} \ar[dd]^{vider}_{0.5}&  \\
            &    & *++[o][F-]{Mort} \ar @(r,u) []_{\emptyset}^1 \\
             &     *++[o][F-]{VV} \ar [ru]^{\emptyset}_{0.5} &  \\
        }
    }
\end{align}

Il est à noter que cet exemple représente un exemple jouet, des choix subjectifs ont été effectués sur la représentation des états, des événements, mais aussi concernant les valeurs attribuées par $expert\_judge_\tau$ qui correspondent au respect d'une norme sociale abstraite (et non statistique). On considère que les formules s'applique sur un exemple fourni sans le remettre en question.

\subsection{Définition des contextes causaux}
On réduit le champ de travail en ne travaillant que sur des contextes causaux partageant le même état de départ $S_0$ et une politique correcte relativement au LTS $<\mathds{S},\tau>$
\begin{itemize}
\item Contexte causal  $\chi_1 = (\pi_1,S_0,\tau)$ avec $\pi_1\in \Pi_\tau$ :


    \begin{minipage}{0.45\textwidth}
        \begin{align*}
            &\pi_1(VO)= \{empoisonner\}, \quad \pi_1(VE)= \{boire\}, \quad \pi_1(mort)=\{\emptyset\} \quad \\ 
            &\pi_1(VV) = \{\emptyset\},\quad \pi_1(VA) = \{boire\},\quad \pi_1(vivant)=\{\emptyset\}
        \end{align*}
        On obtient :
    \begin{itemize}
        \item $\sigma^{\chi_1}= [VO,empoisonner,VE,boire, Mort]$
        \item $\sigma^{\chi_1}_{|state}=[VO,VE,Mort]$
        \item $\sigma^{\chi_1}_{|event}=[empoisonner,boire]$
    \end{itemize}
    \end{minipage}   
    \hfill
    \begin{minipage}{0.45\textwidth}
    \scalebox{0.65}{    \xymatrix@R=1em@C=5em{
            &    *++[o][F-]{Vivant} \ar @(r,u) []_{\emptyset}^1& \\
            &   &    *++[o][F-]{VA} \ar [lu]_{boire}^1  \\
            *++[o][F-]{VO}  \ar [ruu]^{boire}_1 \ar @[red][r]^{empoisonner}_{0.1} \ar[rdd]^{vider}_{0.5} &*++[o][F-]{VE} \ar[ru]^{antidote}_1 \ar@[red][rd]^{boire}_{0.1} \ar[dd]^{vider}_{0.5}&  \\
            &    & *++[o][F-]{Mort} \ar @(r,u) []_{\emptyset}^1 \\
             &     *++[o][F-]{VV} \ar [ru]^{\emptyset}_{0.5} &  \\
        }}

    \end{minipage}


    \item Contexte causal $\chi_2 = (\pi_2,S_0,\tau)$avec $\pi_2\in \Pi_\tau:$

    \begin{minipage}{0.45\textwidth}

    \begin{align*}
        &\pi_2(VO)= \{boire\}, \quad \pi_2(VE)= \{boire\}, \quad \pi_2(mort)=\{\emptyset\} \quad \\ 
        &\pi_2(VV) = \{\emptyset\},\quad \pi_2(VA) = \{boire\},\quad \pi_2(vivant)=\{\emptyset\}
    \end{align*}

    \begin{itemize}
        \item $\sigma^{\chi_2}= [VO,boire, Vivant]$
        \item $\sigma^{\chi_2}_{|state}=[VO,Vivant]$
        \item $\sigma^{\chi_2}_{|event}=[boire]$
    \end{itemize}
    \end{minipage}   
    \hfill
    \begin{minipage}{0.45\textwidth}
    \scalebox{0.8}{    \xymatrix@R=1em@C=5em{
            &    *++[o][F-]{Vivant} \ar @(r,u) []_{\emptyset}^1& \\
            &   &    *++[o][F-]{VA} \ar [lu]_{boire}^1  \\
            *++[o][F-]{VO}  \ar @[red] [ruu]^{boire}_1 \ar [r]^{empoisonner}_{0.1} \ar[rdd]^{vider}_{0.5} &*++[o][F-]{VE} \ar[ru]^{antidote}_1 \ar[rd]^{boire}_{0.1} \ar[dd]^{vider}_{0.5}&  \\
            &    & *++[o][F-]{Mort} \ar @(r,u) []_{\emptyset}^1 \\
             &     *++[o][F-]{VV} \ar [ru]^{\emptyset}_{0.5} &  \\
        }}

    \end{minipage}


    \item Contexte causal $\chi_3 = (\pi_3,S_0,\tau)$ avec $\pi_3\in \Pi_\tau:$

    \begin{minipage}{0.45\textwidth}
        \begin{align*}
            &\pi_3(VO)= \{vider\}, \quad \pi_3(VE)= \{boire\}, \quad \pi_3(mort)=\{\emptyset\} \quad \\ 
            &\pi_3(VV) = \{\emptyset\},\quad \pi_3(VA) = \{boire\},\quad \pi_3(vivant)=\{\emptyset\}
        \end{align*}
        \begin{itemize}
            \item $\sigma^{\chi_3}= [VO,vider, VV,\emptyset,Mort]$ 
            \item $\sigma^{\chi_3}_{|state}=[VO,VV,Mort]$
            \item $\sigma^{\chi_3}_{|event}=[vider,\emptyset]$
        \end{itemize}
        \end{minipage}   
        \hfill
        \begin{minipage}{0.45\textwidth}
        \scalebox{0.8}{    \xymatrix@R=1em@C=5em{
                &    *++[o][F-]{Vivant} \ar @(r,u) []_{\emptyset}^1& \\
                &   &    *++[o][F-]{VA} \ar [lu]_{boire}^1  \\
                *++[o][F-]{VO}  \ar  [ruu]^{boire}_1 \ar [r]^{empoisonner}_{0.1} \ar@[red][rdd]^{vider}_{0.5} &*++[o][F-]{VE} \ar[ru]^{antidote}_1 \ar[rd]^{boire}_{0.1} \ar[dd]^{vider}_{0.5}&  \\
                &    & *++[o][F-]{Mort} \ar @(r,u) []_{\emptyset}^1 \\
                 &     *++[o][F-]{VV} \ar @[red][ru]^{\emptyset}_{0.5} &  \\
            }}
    
        \end{minipage}

    \item Contexte causal 
    $\chi_4 = (\pi_4,S_0, \tau)$ avec $\pi_4\in \Pi_\tau:$

    \begin{minipage}{0.45\textwidth}
    \begin{align*}
        &\pi_4(VO)= \{empoisonner\}, \quad \pi_4(VE)= \{vider\}, \quad \pi_4(mort)=\{\emptyset\} \quad \\ 
        &\pi_4(VV) = \{\emptyset\},\quad \pi_4(VA) = \{boire\},\quad \pi_4(vivant)=\{\emptyset\}
    \end{align*}
    \begin{itemize}
        \item $\sigma^{\chi_4}= [VO,empoisonner,VE,vider,VV,\emptyset,Mort]$
        \item $\sigma^{\chi_4}_{|state}=[VO,VE,VV,Mort]$
        \item $\sigma^{\chi_4}_{|event}=[empoisonner,vider,\emptyset]$
    \end{itemize}        
    \end{minipage}
    \hfill
    \begin{minipage}{0.45\textwidth}
        \scalebox{0.65}{    \xymatrix@R=1em@C=5em{
                &    *++[o][F-]{Vivant} \ar @(r,u) []_{\emptyset}^1& \\
                &   &    *++[o][F-]{VA} \ar [lu]_{boire}^1  \\
                *++[o][F-]{VO}  \ar  [ruu]^{boire}_1 \ar@[red] [r]^{empoisonner}_{0.1} \ar[rdd]^{vider}_{0.5} &*++[o][F-]{VE} \ar[ru]^{antidote}_1 \ar[rd]^{boire}_{0.1} \ar@[red][dd]^{vider}_{0.5}&  \\
                &    & *++[o][F-]{Mort} \ar @(r,u) []_{\emptyset}^1 \\
                 &     *++[o][F-]{VV} \ar@[red][ru]^{\emptyset}_{0.5} &  \\
            }}
    \end{minipage}

    

    \item Contexte causal $\chi_5 = (\pi_5,S_0,\tau)$ avec $\pi_5\in \Pi_\tau:$


     \begin{minipage}{0.3\textwidth}
        \begin{align*}
            &\pi_5(VO)= \{empoisonner\}, \quad \pi_5(VE)= \{antidote\}, \quad \pi_5(mort)=\{\emptyset\} \quad \\ 
            &\pi_5(VV) = \{\emptyset\},\quad \pi_5(VA) = \{boire\},\quad \pi_5(vivant)=\{\emptyset\}
        \end{align*}
         \begin{itemize}
             \item $\sigma^{\chi_5}= [VO,empoisonner,VE,antidote,VA,boire,Vivant]$ 
             \item  $\sigma^{\chi_5}_{|state}=[VO,VE,VA,Vivant]$
             \item $\sigma^{\chi_5}_{|event}=[empoisonner,antidote,boire]$
         \end{itemize}         
     \end{minipage}
     \begin{minipage}{0.45\textwidth}
        \scalebox{0.6}{    \xymatrix@R=1em@C=5em{
                &    *++[o][F-]{Vivant} \ar @(r,u) []_{\emptyset}^1& \\
                &   &    *++[o][F-]{VA} \ar @[red][lu]_{boire}^1  \\
                *++[o][F-]{VO}  \ar  [ruu]^{boire}_1 \ar@[red] [r]^{empoisonner}_{0.1} \ar[rdd]^{vider}_{0.5} &*++[o][F-]{VE} \ar@[red][ru]^{antidote}_1 \ar[rd]^{boire}_{0.1} \ar[dd]^{vider}_{0.5}&  \\
                &    & *++[o][F-]{Mort} \ar @(r,u) []_{\emptyset}^1 \\
                 &     *++[o][F-]{VV} \ar[ru]^{\emptyset}_{0.5} &  \\
            }}
     \end{minipage}
\end{itemize}

\subsection{Définition des politiques optimales}

Politique optimale sans planification :
\begin{align*}
    & \pi^{opt1}_\tau(VO)=\{boire\}, \quad \pi^{opt1}_\tau(VE)=\{antidote\}, \quad \pi^{opt1}_\tau(VV)=\{\emptyset\}, \quad \pi^{opt1}_\tau(Mort)=\{\emptyset\} \\
    & \pi^{opt1}_\tau(VA)=\{boire\}, \quad \pi^{opt1}_\tau(Vivant)=\{\emptyset\} \quad
\end{align*}
Politique optimale avec planification :
\begin{align}\label{pol_opti}
    & \pi^{opt2}_\tau(VO)=\{boire\}, \quad \pi^{opt2}_\tau(VE)=\{antidote\}, \quad \pi^{opt2}_\tau(VV)=\{\emptyset\}, \quad \pi^{opt2}_\tau(Mort)=\{\emptyset\} \\
    & \pi^{opt2}_\tau(VA)=\{boire\}, \quad \pi^{opt2}_\tau(Vivant)=\{\emptyset\} \quad
\end{align}
On note $\forall S\in \mathds{S},\quad \pi^{opt1}_\tau(S) =\pi^{opt2}_\tau(S)$. On notera par la suite simplement $\pi^{opt}_\tau$.
\begin{figure}[h!]
    \centering
    \[
        \xymatrix@R=1em@C=5em{
            &    *++[o][F-]{Vivant} \ar@[red]  @(r,u) []_{\emptyset}^1& \\
            &   &    *++[o][F-]{VA} \ar@[red]  [lu]_{boire}^1  \\
            *++[o][F-]{VO}  \ar @[red] [ruu]^{boire}_1 \ar [r]^{empoisonner}_{0.1} \ar[rdd]^{vider}_{0.5} &*++[o][F-]{VE} \ar@[red] [ru]^{antidote}_1 \ar[rd]^{boire}_{0.1} \ar[dd]^{vider}_{0.5}&  \\
            &    & *++[o][F-]{Mort} \ar@[red]  @(r,u) []_{\emptyset}^1 \\
             &     *++[o][F-]{VV} \ar@[red] [ru]^{\emptyset}_{0.5} &  \\
        }    
    \]
    \captionsetup{width=0.8\textwidth}
    \caption{Politique optimale: les flèches rouges indiquent dans chaque état l'événement choisi par la politique optimale}
    \label{fig:placeholder}
\end{figure}

\subsection{Expérimentations}

En fonction de la configuration choisie (extraction d'une politique optimale, exploitation directe, évaluation de la normalité globale, sur les événements avec alternatives uniquement) et des hyper paramètres ($\lambda_1,\lambda_2, \gamma$...) on obtient différents résultats. Une connaissance à priori des contraintes que l'on souhaite appliquer lors de l'évaluation permet de guider le classement. 

Dans les exemples suivant on considère, si ce n'est pas précisé, que le contexte causal observé est $\chi_1$, on fixe  $\lambda_1 =\lambda_2 =1$ et un coût maximal $\delta=2$ (avec cette valeur la validité ne dépends que de la différence entre les deux états finaux des traces de l'observation et du contrefactuel).

\paragraph{Priorité à des contrefactuels avec des scores d'événements consistants :}
On choisit pour la fonction de vraisemblance la formule \ref{exploit_direct} avec comme opérateur d'agrégation $agg_{geo}$ et on agrège les scores de tous les événements des scénarios. On souhaite dans cette configuration pénaliser des traces avec une variance de score importante.

On obtient les résultats suivants:

\begin{table}[h]
\centering
\begin{tabular}{|c|c|c|c|c|c|c|}
\hline
CF & classement & $d_{\text{obs}}$ & $v$ & $d_{\text{obs}} + v$ & trace & valide \\
\hline\hline
1 & $\chi_2$ & 0.67 & 0.0 & 0.67 &  $[VO,boire, Vivant]$ & \checkmark \\
\hline
3 & $\chi_3$ & 0.33 & 0.5 & 0.83 & $[VO,vider, VV,\emptyset,Mort]$ & \texttimes \\
4 & $\chi_4$ & 0.25 & 0.71 & 0.96 & $[VO,empoisonner,VE,vider,VV,\emptyset,Mort]$ &\texttimes \\
2 & $\chi_5$ & 0.50 & 0.54 & 1.04 & $[VO,empoisonner,VE,antidote,VA,boire,Vivant]$ & \checkmark \\
\hline
\end{tabular}
\caption{Résultats pour la priorité à la normativité}
\label{tab:prio_normativite}
\end{table}

On observe dans la table \ref{tab:prio_normativite} que le contexte contrefactuel avec le coût minimal est le contexte $\chi2$ sa distance à l'observation est de 0.67 et sa vraisemblance de 0 (l'unique événement $boire$ a un score de 1). Ce contrefactuel est suivi du contexte $\chi_5$, plus proche de l'observation (coût de 0.5) mais le scénario associé contient l'événement $empoisonner$ qui fait augmenter le coût. Les deux premiers contrefactuels $\chi_2$ et $\chi_5$ sont valides.



\paragraph{Interdiction de dépasser un seuil}

On utilise l'opérateur d'agrégation avec un seuil fixé à $\lambda = 0.2$ pour pénaliser les contrefactuels dont le scénario contient un empoisonnement. On fixe aussi un coût maximal de 0.5 pour le coût de vraisemblance. 

\begin{table}[h]
\centering
\begin{tabular}{|c|c|c|c|c|c|c|}
\hline
CF & classement & $d_{\text{obs}}$ & $v$ & $d_{\text{obs}} + v$ &  trace & valide \\
\hline\hline
2 & $\chi_2$ & 0.67 & 0.0 & 0.67 & $[VO,boire, Vivant]$ &  \checkmark \\
\hline
1 & $\chi_3$ & 0.33 & 0.0 & 0.33 & $[VO,vider, VV,\emptyset,Mort]$ & \texttimes \\
3 & $\chi_4$ & 0.25 & 1.0 & 1.25 & $[VO,empoisonner,VE,vider,VV,\emptyset,Mort]$ & \texttimes \\
4 & $\chi_5$ & 0.50 & 1.0 & 1.50 & $[VO,empoisonner,VE,antidote,VA,boire,Vivant]$ &  \texttimes \\
\hline
\end{tabular}
\caption{Résultats pour la priorité à la normativité}
\label{tab:res_seuil}
\end{table}

Dans cette configuration uniquement le contexte causal $\chi_2$ est valide. En rajoutant une contrainte sur le coût de vraisemblance, on exclut automatiquement les contrefactuels qui possèdent un événement inférieur au seuil, c'est ce qu'on observe pour $\chi_5$. 

\paragraph{Normativité en début de scénario}
On utilise maintenant l'opérateur $agg_\gamma$, avec $\gamma = 0.25$, qui avantage les scénarios qui ont des scores élévés sur les événements en début de trace en considérant peu ceux de la fin de la trace.

\begin{table}[h]
\centering
\begin{tabular}{|c|c|c|c|c|c|c|}
\hline
CF& classement & $d_{\text{obs}}$ & $v$ & $d_{\text{obs}} + v$ & trace & valide \\
\hline\hline
1 & $\chi_2$ & 0.67 & 0.0 & 0.67 & $[VO,boire, Vivant]$ & \checkmark \\
\hline
2 & $\chi_3$ & 0.33 & 0.5 & 0.83 & $[VO,vider, VV,\emptyset,Mort]$ & \texttimes \\
3 & $\chi_4$ & 0.25 & 0.80 & 1.05 & $[VO,empoisonner,VE,vider,VV,\emptyset,Mort]$ & \texttimes \\
4 & $\chi_5$ & 0.50 & 0.69 & 1.19 & $[VO,empoisonner,VE,antidote,VA,boire,Vivant]$ & \checkmark \\
\hline
\end{tabular}
\caption{Priorité à la normativité en début de scénario}
\label{tab:res_gamma}
\end{table}

On observe avec la table \ref{tab:res_gamma} que le contrefactuel avec le plus faible coût correspond à $\chi_2$ qui a le premier événement avec un score le plus élevé ($boire$ =1) suivit de $\chi_3$ qui a un score pour son premier événement de 0.5. Les deux derniers contextes causaux possède en premier événement $empoisonner$ ce qui vient fortement augmenter le coût. 


\paragraph{Normativité en fin de scénario}

On utilise l'opérateur $gamma_{late}$ avec $\gamma =0.25$. On priorise la normativité en fin de trace, couplé au coût sur la distance à l'observation on cherche des contrefactuels qui ne diffèrent de l'observation qu'en fin de trace mais qui reste le plus proche possible au début.

\begin{table}[h]
\centering
\begin{tabular}{|c|c|c|c|c|c|c|}
\hline
CF & classement & $d_{\text{obs}}$ & $v$ & $d_{\text{obs}} + v$ & trace & valide \\
\hline\hline
2 & $\chi_2$ & 0.67 & 0.0 & 0.67 & $[VO,boire, Vivant]$ & \checkmark \\
4 &$\chi_3$ & 0.33 & 0.5 & 0.83 & $[VO,vider, VV,\emptyset,Mort]$ &  \texttimes \\
3 & $\chi_4$ & 0.25 & 0.52 & 0.77 & $[VO,empoisonner,VE,vider,VV,\emptyset,Mort]$ & \texttimes \\
\hline
1 & $\chi_5$ & 0.50 & 0.04 & 0.54 & $[VO,empoisonner,VE,antidote,VA,boire,Vivant]$ &  \checkmark \\
\hline\hline
\end{tabular}
\caption{Priorité à la normativité en fin de scénario}
\label{tab:res_gamma_late}
\end{table}

La table des résultats \ref{tab:res_gamma_late} montre que dans cette configuration, c'est $\chi_5$ qui a le coût le plus faible, malgré l'empoisonnement, on s'éloigne de l'observation le plus tard possible.

\paragraph{Normativité du contrefactuel avec distance à une politique optimale}
On choisit pour cette configuration de calculer la distance à la politique optimale définie par \ref{pol_opti}. Le choix est fait de ne considérer que les événements pour lesquelles une alternative existe.

\begin{table}[h]
\centering
\begin{tabular}{|c|c|c|c|c|c|c|}
\hline
CF& classement & $d_{\text{obs}}$ & $v$ & $d_{\text{obs}} + v$ & trace & valide \\
\hline
1 & $\chi_2$ & 0.67 & 0.0 & 0.67 & $[VO,boire, Vivant]$ & \checkmark \\
\hline
4 & $\chi_3$ & 0.33 & 1.0 & 1.33 & $[VO,vider, VV,\emptyset,Mort]$ & \texttimes \\
3 & $\chi_4$ & 0.25 & 1.0 & 1.25 & $[VO,empoisonner,VE,vider,VV,\emptyset,Mort]$ & \texttimes \\
2 & $\chi_5$ & 0.50 & 0.5 & 1.0 & $[VO,empoisonner,VE,antidote,VA,boire,Vivant]$ & \checkmark \\
\hline
\end{tabular}
\caption{Résultats distance à la politique optimale.}
\label{tab:dist_pol}
\end{table}

Les résultats présents dans la table \ref{tab:dist_pol} indiquent qu'avec cette configuration, le contexte causal avec le plus faible coût est le contexte $\chi_2$ suivi de $\chi_5$ les deux contrefactuels valides. Concernant la fonction de vraisemblance : dans la trace associée à $\chi_2$ un seul événement est présent et un choix est possible, le choix fait est l'événement $boire$ qui correspond à l'événement choisi par la politique optimale donc le coût associé est de 0. Pour $\chi_5$, deux événements avec alternative sont présents dans la trace $empoisonner$ qui n'est pas l'événement indiqué par la politique optimale, et $antidote$ qui est similaire à la politique optimale. Sur les deux événements un seul est égal à la politique optimale donc un coût de 0.5 est appliqué. Pour $\chi_3$ et $\chi_4$ tous les choix possibles différents de la politique optimale un coût d'un leur est attribué.





